\chapter[Sermon 8. The Second Part of the Epistle to the Ephesians Where Is Spoken of Penance and of the Nicolaites.]{The Second Part of the Epistle to the Ephesians Where Is Spoken of Penance and of the Nicolaites.}
\label{sermon:008}
\markright{Sermon 8. The Second Part of the Epistle to the Ephesians Where ...}

Remember therefore from whence you have fallen, and repent, and do your first works. Or else I will come to you shortly, and will remove your lampstand out of its place, except you repent. But you have this, because you hate the deeds of the Nicolaitanes, which deeds I hate also. Let him who has ears, hear what the Spirit says to the congregations. To him that overcomes will I give to eat of the tree of life, which is in the midst of the Paradise of God.

\index{Repentance}The accusations our Savior Christ uses against his servants who sin do not, I am certain, aim to cause those overwhelmed with reproaches to be ashamed, despair, and perish; but rather that they should be restored and live. Therefore, the Lord Jesus immediately adds an exhortation to repentance, that they may be saved, and also describes the true and lawful penalty.

For we heard what he rebuked in the congregations of the Ephesians; let us hear now what the Lord requires of them, and how he seeks to have the error corrected truly by repentance, to which he exhorts. For he has said that the Lord strikes and heals; chiefly in this case. Which doctrine surely is proper and perpetual to the church of Christ.

He speaks primarily of three things in this matter, in his counsel or exhortation to amendment. First, he reminds them of from what they have fallen: that is to say, with how great love they had burned previously, and how they could now be deceived. He shows how fortunate and blessed a state they had maintained until now, and how unfortunate and shameful their current state is. For acknowledging the transgression is the beginning of repentance, if, illuminated by faith, we consider what benefits we have lost and in what misery we now find ourselves. He who believes he has lost nothing will not be moved to make any search or inquiry; he who thinks he has fallen from no happiness will not consider how he may be restored. Therefore, in the matter of a godly life, acknowledgment and confession of sins must come first, by which we may lament before God our poverty and misery. Indeed, they do not fall from happiness who were never in any happiness. Therefore, we say that holy men may fall, and also be restored by repentance. Then, after the acknowledgment of our failure, follows repentance: that is, the turning again of our mind, so that we do not proceed forever like madmen and fools, walking in the way of vanity and unrighteousness. For repentance is a conversion and change, whereby we are turned away from evil and converted to good, returning to our senses and working righteousness.

Finally, what follows explains this repentance. Repeat your former deeds: be fervent again in love, perform the works of faith, which are fruits worthy of repentance. See, there is no need for new laws or lengthy debates concerning amendment. Briefly, repeat your former deeds, not of the flesh, but those which you began when you first received the Gospel and were born again in Christ. Ultimately, this is the true reformation: to do the former works of God, not the latter, which the error of the world has devised. This truly is, and will be, the true repentance: namely, the acknowledgment of sin, a turning to God and to good, and a turning away from the Devil and from evil, and the working of the first righteousness, the righteousness of faith in Christ. There are many and varied discussions about repentance and its parts, concerning the contrition of the heart, the confession of the mouth, and the satisfaction of that work. But just as none is briefer than that of Jesus Christ, so truly is there none better or more certain.

\index{Arethas of Caesarea}\index{Mahomet}\index{Turks}To these his exhortations and godly cautions, he adds most grievous threatenings, if a great danger could not be avoided. And he speaks few words indeed, but understands a great evil that cannot be spoken or declared. “Unless you repent, I will remove your lampstand from its place.” The lampstand, as the Lord himself has explained, is the church. It stands in its place while it leans toward Christ, and is preserved by Christ so long as the preaching of the truth is maintained within it. It is indeed the church of Christ. She is removed from her place when she is without the preaching of the truth, and no longer leans upon Christ, nor is defended by him, but is forsaken, and is no longer indeed the church of Christ. This is done by Christ himself, through his just judgment, when our unthankfulness and a life that cannot repent drives God to depart from us, to abandon us to our error and darkness, and to leave us to deceivable men, and so on. This understanding Arethas acknowledges, who says that to remove the church is when it is left bare and destitute of God’s grace, because of which nakedness it wavers in doubtful perplexity and in storms cast upon it by wicked men. And truly we see how at this day the church of the Ephesians is removed from its place, and no longer enjoys the wholesome doctrine of Christ, nor stands upon the wholesome rock Jesus Christ, but is oppressed by the pestiferous doctrine, or rather madness, of Muhammad, and lies in sorrow under the feet of the Turks. We see at this day in Germany (and it is the greater pity) many lampstands removed from their place, not without the great triumph of Satan and the loss of souls. Moreover, it is also to be observed in this threatening that, without terror, he says, “I will come to you shortly.” For it is a phrase of speech. For we also say, “I will come to you by and by,” that is, I will come to avenge and punish, and perhaps sooner than you look for. Most certainly, whenever I happen to come, I will take punishment of you. Let no man therefore think to escape unpunished in a life that cannot repent.

Again, where the Lord repeats, “except you repent,” he plainly testifies that the bosom of God’s mercy and clemency is readily open, if we do penance, however we have offended him before. In the meantime, we learn here openly and most certainly that we can by no counsels or consultations, by no armies nor policies, prevail against our perils unless we repent. Therefore, unless we will have our churches subverted and given over to be seduced and destroyed by the devil and his seducers, let us repent in time and receive again the first love.

Again he commends the singular virtue in this congregation, especially for that they have hated the doings of the Nicolaitanes, which God himself also hates. Let us mark every word here. He does not say you did slyly or avoid and abstain, but that you have hated. The force of hatred is great, moving one to persecute, that you hate.

Moreover, he does not say, “You have hated the Nicolaitans,” but “the works of the Nicolaitans.” For we ought to hate no man’s person in itself, but the vice in the man, so that when we forsake it, we should love the man with all our heart. It must needs be a great evil, which God himself confesses that he hates. Here all congregations shall understand that they ought also by all means to hate the heresy and abomination of the Nicolaitans. Albeit that at this day the name be extinguished, yet the heresy and abomination of the Nicolaitans remains.

\index{Augustine, Saint}\index{Irenaeus}\index{Tertullian}\index{Eusebius of Caesarea}\index{Epiphanius of Salamis}This Nicolas was from Antioch, one of the seven deacons, mentioned in chapter 6 of the Acts. He is said to have revolted from the purity of faith, as Judas did. And where he was previously a Gentile (for it is said that he was a proselyte), he returned in some respects to gentility, like a dog to its vomit. The Nicolaitans are also Gnostics and associates of Carpocrates, filthy and most wicked people. Clement somewhat excuses Nicolas in Eusebius in book 3, chapter 29 of the \textit{Ecclesiastical History}. But that excuse seems insufficient or just, since all the ancients accuse him with one voice, and namely, the judgment of God in this present and in the following epistles. Irenaeus condemns him from this same place, in book 1 against the Valentinians, chapter 27 and following. Tertullian, at the end of Prescribed Heretics, touches cleverly on the facts of the Nicolaitans and detests them. Nevertheless, he does not explain them, but passes them over. And I do not know how cleverly Epiphanius has uttered and declared the wicked and abominable acts, neither to be thought nor told, and most beastly filthiness, such as has not been heard of in the heresy, 25, 26, 27, and 31 and following. Philastrius and Augustine also have touched on the Nicolaitans, either of them in their lists of heresies. Shameful details will not allow me to recite. It is enough if we know what the Lord himself explained in the epistle to Pergamum, calling the doctrine of the Nicolaitans the doctrine of Balaam the false prophet. But who does not know what counsel he gave to King Balaac of Moab and Midian, and how he prostituted fair women to the young men of Israel, by whose acquaintance they both defiled themselves with fornication and ate also of meats offered up to idols, becoming partakers of Baal-peor. Let him who will read Josephus in book 4 of Antiquities, chapter 6. And without doubt, the sacrifices of the Nicolaitans seem to differ nothing from the secrets of Priapus, or Gerecinthia, or the mother of God, and the nightly service of Bacchus. Irenaeus openly signifies that the Carpocratites, who are also called Gnostics, did not abhor images, but painted and fashioned into images Jesus and Paul, along with images of certain philosophers. And that the image of Jesus, as they say, was made expressly by Pilate, who commanded that the face of Jesus be painted vividly. But however that was, it is certain that the actions of the Nicolaitans were in ill repute because of their fornications and adulteries. And that the Nicolaitans abstained not from images, nor from meats offered to idols. Against this error, Paul also wrote many things.

Hereof let us learn to abhor and fly fornication, and never to think of restoring brothels or other places of whoredom. Fie for shame. Let us learn hereby to keep virginity, single life, and lawful marriages, and avoid those called Nicolaitans. Let us learn hereby to keep ourselves well from idols, idolatry, and from all strange kinds of worshipings. All those God hates.

And with an acclamation he pierced the ears of all, moving all to attentiveness and holy obedience. And he applied this doctrine to all times and to all congregations throughout the world. He used his accustomed speech, repeated so often in the Gospel: “He that hath ears to hear, let him hear.” It is not in our strength to hear and obey God. For God prepares our ears, and with his grace frames and draws our hearts. And let those to whom the grace of God is given, beware lest through their negligence, vanity, and lightness they lose it not. Let them show such diligence as God in his word requires and prescribes. They that do have ears to hear. He says therefore, take heed to what God now speaks, and whose hearts he now stirs and moves, that you lose not this grace through your negligence. Be diligent, attentive, and circumspect, stirring up within yourselves the gift of God.

Now also he urges diligence by authority. The Spirit of God speaks and reveals these things, not the spirit of men or of error, for God speaks by his scripture, which is recognized as the spirit both of the Father and of the Son. Moreover, he applies all things and every thing to all congregations, where he says, “What the Spirit says to the congregations,” not to the congregation. It is now manifest, and beyond all controversy, that those seven churches represent a figure of all churches throughout the whole world, and that all are instructed in those seven.

Moreover, to ensure that nothing is lacking in the just exhortation to repentance, faith, and diligence, he adds a most ample promise and employs allegorical language, so that it might have the greater grace. To those who overcome, he promises to give the fruit of the tree of life, planted in the paradise of God. He alludes to the second chapter of Genesis and transfers the meaning from earthly things to heavenly ones. The paradise of God (which some understand to be the church) is that everlasting blessedness and felicity of which the Lord spoke to the thief, saying, “This day you will be with me in paradise.” Herein is the tree of life, Christ, communicating his eternal life to us, which we enjoy and experience while being conveyed into heaven by him and living with him. Finally, this is that ambrosia or godly drink which the heavenly Father gives us to drink. But this great and wonderful good does not happen to everyone, but only to him who overcomes. For Adam did not overcome; he had died in defeat. Therefore, if we overcome the flesh, the Devil, and the world, and that through Christ, we shall also live in the world to come with Christ.

The Complutensian Bible has it, which is in the midst of the Paradise of my God. And Arethas explains it, and says: Let no man be offended by this. All humble things agree to the dispensation of the incarnation, which was made for our cause, since he himself in the Gospel says: I ascend unto my Father, and your Father, to my God, and to your God. \&c.

And thus far, concerning the Epistle of Jesus Christ by John to the Ephesians, and what benefit our churches, and each of us, may receive from it. May the Lord enlighten the eyes of our minds.
